% =====================================================================
% paper_jsac_companion.tex
% A companion explainer for Robin, walking him through paper_jsac.tex.
% Not for submission. Casual voice. Single column. Lots of figures.
% Targets: the comms-side jargon (WMMSE, KKT, ZF generalisation, scene
% dictionary, tier system, regime story, Cauchy failure) that the main
% paper takes for granted but Robin does not yet have a clean mental
% model for.
% =====================================================================

\documentclass[11pt]{article}

\usepackage[a4paper,margin=2.4cm]{geometry}
\usepackage[utf8]{inputenc}
\usepackage[T1]{fontenc}
\usepackage{lmodern}
\usepackage{microtype}
\usepackage{parskip}

\usepackage{amsmath,amssymb,bm,mathtools}
\usepackage{graphicx}
\graphicspath{%
  {../../JSAC/code/experiments/rank_check/}%
  {../../JSAC/code/experiments/cauchy_tightness/}%
  {../../JSAC/code/experiments/csi_calibration/}%
  {../../JSAC/code/experiments/tier_c_decision/}%
  {../../JSAC/code/experiments/plaza_run/figures/}%
}
\usepackage{booktabs}
\usepackage{caption}
\usepackage[hidelinks]{hyperref}
\usepackage{xcolor}

\usepackage{enumitem}

% Sidebar for "the part that confused you"
\usepackage[framemethod=TikZ]{mdframed}
\newenvironment{nb}[1][]{%
  \begin{mdframed}[backgroundcolor=yellow!7,linecolor=yellow!50!black,
    linewidth=0.4pt,roundcorner=2pt,innertopmargin=8pt,
    innerbottommargin=8pt,innerleftmargin=10pt,innerrightmargin=10pt,
    frametitle={\textbf{#1}},frametitlebackgroundcolor=yellow!20]%
}{\end{mdframed}}
\newenvironment{defs}[1][]{%
  \begin{mdframed}[backgroundcolor=blue!4,linecolor=blue!40!black,
    linewidth=0.4pt,roundcorner=2pt,innertopmargin=8pt,
    innerbottommargin=8pt,innerleftmargin=10pt,innerrightmargin=10pt,
    frametitle={\textbf{#1}},frametitlebackgroundcolor=blue!10]%
}{\end{mdframed}}

% math macros (kept aligned with paper_jsac.tex so you can flip between)
\newcommand{\khat}{\hat{\bm{k}}}
\newcommand{\nhat}{\hat{\bm{n}}}
\newcommand{\rr}{\mathbf{r}}
\newcommand{\xx}{\mathbf{x}}
\newcommand{\hh}{\mathbf{h}}
\newcommand{\ww}{\mathbf{w}}
\newcommand{\gvec}{\mathbf{g}}
\newcommand{\Pabs}{P_{\mathrm{abs}}}
\newcommand{\Sab}{S_{\mathrm{ab}}}
\newcommand{\Aab}{A_{\mathrm{ab}}}
\newcommand{\Qmat}{\mathbf{Q}}
\newcommand{\Mmat}{\mathbf{M}}
\newcommand{\Jmat}{\mathbf{J}}
\newcommand{\Phimat}{\mathbf{\Phi}}

\title{\textbf{A companion to the JSAC paper, written for Robin}\\[2pt]
\large{Big picture, comms jargon decoded, and what the experiments are actually saying}}
\author{Self-reading edition. Not for submission.}
\date{2026-05-07}

\begin{document}
\maketitle

\section*{Read this first}

You wrote \texttt{paper\_jsac.tex} (and \texttt{paper\_v2.tex} before it,
and the spine doc, and three rounds of brainstorm). All of it sits on
\texttt{ARCHEOLOGY.md} and \texttt{paper\_spine.md}. The paper itself is
correct and self-contained, but you are not the audience: a JSAC reviewer
who already lives in multi-user MIMO is. So when you re-read it, half
the sentences sit on terminology that does not yet have a clean home in
your head, and you cannot tell whether a paragraph is doing real work or
just paying its taxes.

This document is a friend reading the paper with you. It assumes you
already know everything in:
\begin{itemize}[leftmargin=1.4em,itemsep=2pt]
  \item \texttt{theory/summary\_paper.tex} (the monograph summary): the
  local absorption law, Cauchy/projected area, the exposure operator
  $\Qmat$, the Fresnel/depth factorisation.
  \item \texttt{papers/TAP\_paper/paper.tex}: closed-form whole-body
  dosimetry, $T_0$, $\Aab$, ambient occlusion, layered fat below 6\,GHz.
\end{itemize}
It will not re-explain those. It will explain the \textbf{communications}
side, the \textbf{architecture}, and the \textbf{regime story} the JSAC
paper actually rests on, plus the half-dozen places you got lost.

Section 1 is the $30$-second pitch in your own voice.
Section 2 is a defs of the comms terms you have been seeing without a
dictionary. Section 3 walks the precoder math at the level your
backend brain wants. Section 4 explains the closed-loop architecture
(why the dictionary, translation phasor, CSI calibration, tiers all
need to be there). Section 5 tells you what each plaza-run experiment
is actually doing, with the figure inline so you do not have to flip.
Section 6 is the Cauchy story (it is more interesting than you remember).
Section 7 is what is genuinely new in this paper vs. what you already
knew from earlier work, and what the next-paper hooks are.

Every figure is dropped in here so you can read it without bouncing.

\bigskip\hrule\bigskip

\section{The 30-second pitch in your own voice}

\textbf{One sentence:} every body in the cell becomes a $64\!\times\!64$
Hermitian matrix $\Qmat^{(u)}$ inside the base station's precoder, the
multi-user QCQP under per-body absorbed-power caps has a closed form
that contains zero-forcing as its rank-one limit, and the plaza
experiment shows the cap is slack by four orders of magnitude so the
practical value is chronic-dose stratification by tier, not
instantaneous compliance.

\textbf{Three-paragraph version:}

\textit{The blocker.} Brussels caps cumulative outdoor field at
$14.57\,$V/m. Italy caps at 6 to $15\,$V/m. These are reference levels:
free-space proxies for actual tissue absorption, and they are 10 to
$30\times$ tighter than ICNIRP's basic restriction. Operators
deploying mmWave downlink in those jurisdictions are not propagation-limited,
they are RL-limited. The cap is conservative because it cannot see what
is actually happening on the body. Your machinery can.

\textit{The move.} Treat each body as a digital twin maintained inside
the BS control loop. Each body contributes a Hermitian PSD operator
$\Qmat^{(u)} \in \mathbb{C}^{M \times M}$ such that absorbed power is
$\xx^H \Qmat^{(u)} \xx$. The BS already does ray tracing on its scene
model to plan beams. The same path set that gives the channel
$\hh_k$ to the served user gives every $\Qmat^{(u)}$ on every body
through your TAP-paper machinery. So the twin is computed from
quantities the network already maintains. Stack budgets across served
users, cooperating non-users, and ISAC-detected bystanders, and solve
the multi-user sum-rate QCQP. The KKT closed form is
$\ww_k^\star = (\Qmat_{\mathrm{tot}} + \nu \mathbf{I})^{-1} \gvec_k$.
That is the same shape as regularised zero-forcing. ZF is its
rank-one degenerate limit. Drop-in, not a new solver.

\textit{The reality check.} On the Brussels Grand Place geometry with
50 SMPL-X bodies at $26\,$GHz on an $8\!\times\!8$ panel and 25 served
users, the cap sits four orders of magnitude away from the actual
absorbed power. The five precoders you compare (MRT, ZF, worst-case
back-off, multi-body ECBF, pose-known oracle) collapse to one operating
point on the budget metric. So the paper pivots: it characterises the
\emph{regime} (when does the cap bind, when does it not), shows that
chronic-dose ECDFs separate by tier even when instantaneous compliance
is easy ($\sim 1.5\times$ between served and cooperating), and walks
back the Cauchy bound from a universal upper bound to a conditional
one (it is loose by 21 to $27\,$dB on the steering direction and
breached by half of random precoders).

That is the paper. Now everything else is filling in why each piece is
there.

\section{Glossary: the comms side, in your language}

You know what a precoder is. You know MRT. You know what "illuminate a
phantom and compute Sab" means. The paper assumes a half-dozen more
things that you have never had a reason to internalise. Here they are.

\begin{defs}[Precoding terms cheat sheet]
\textbf{Precoder $\xx$ or $\ww_k$.} A complex vector in $\mathbb{C}^M$
that the BS multiplies the data symbol by, then transmits one component
per antenna element. So if $M=64$, the precoder is 64 complex numbers,
one phase-and-amplitude knob per element. The radiated field anywhere
in space is linear in $\xx$. Every dosimetric quantity that depends on
the field is quadratic in $\xx$. That is why $\Qmat$ shows up: power
is $|E|^2$, $E$ is linear in $\xx$, so power is $\xx^H \Qmat \xx$ for
some Hermitian $\Qmat$. You wrote that operator. The whole paper is
about exploiting that quadratic form.

\textbf{Channel $\hh_k$.} For served user $k$, the BS-to-UE channel
vector. The received baseband signal at user $k$ is $y_k = \hh_k^T \xx
+ n$. The channel comes from the same ray tracer that gives you the
path set used to build $\Qmat^{(u)}$. Your TAP paper computes
$\Sab(\rr)$ from a path set; the network computes $\hh_k$ from the
same path set evaluated at the UE position. Same upstream, two
downstream uses.

\textbf{MRT (maximum ratio transmission).} The simplest single-user
precoder: $\xx \propto \hh_k^*$. Phases align constructively at the UE.
Maximises SNR. Ignores everyone else. You know this.

\textbf{ZF (zero-forcing).} Multi-user. Stack $K$ user channels as the
rows of $\mathbf{H} \in \mathbb{C}^{K \times M}$. The ZF precoder for
user $k$ is the $k$-th column of $\mathbf{H}^H (\mathbf{H} \mathbf{H}^H)^{-1}$,
which projects the precoder for user $k$ onto the null space of all
other users' channels. So user $k'$ sees $\hh_{k'}^T \ww_k = 0$ for
$k' \neq k$. Hard nulls, no inter-user interference, but it costs DoF:
each null eats one degree of freedom. With $K=M$, you exhaust the
array.

\textbf{Regularised ZF.} Replace the inverse with
$(\mathbf{H} \mathbf{H}^H + \nu \mathbf{I})^{-1}$. Soft nulls, trades
a bit of interference for not blowing up when $\mathbf{H} \mathbf{H}^H$
is near-singular. Peel-Hochwald-Swindlehurst 2005.

\textbf{WMMSE (weighted minimum mean square error).} The standard
trick for converting a non-convex multi-user sum-rate maximisation
into something solvable. The sum-rate $\sum_k \log(1 + \mathrm{SINR}_k)$
is non-convex in $\{\ww_k\}$ jointly. WMMSE replaces it with a quadratic
that has the same fixed points: $\sum_k [\log u_k - u_k e_k(\ww_k)]$,
where $e_k$ is the per-user MSE and $u_k$ is a per-user weight. You
alternate three steps: optimise $u_k$, optimise the receive equaliser
$v_k$, optimise $\ww_k$. The third step is quadratic in $\ww_k$, so it
plays nicely with quadratic exposure constraints. Christensen 2008,
Shi 2011. \textbf{This is why the closed form exists.} Without WMMSE,
the multi-user sum-rate problem under exposure caps does not collapse
to one matrix inverse per user.

\textbf{KKT / Karush-Kuhn-Tucker.} The first-order optimality conditions
for a constrained optimisation. For a Lagrangian
$\mathcal{L} = f + \sum_i \lambda_i g_i$ with constraints $g_i \leq 0$,
KKT says: stationarity ($\nabla f + \sum \lambda_i \nabla g_i = 0$),
primal feasibility ($g_i \leq 0$), dual feasibility ($\lambda_i \geq 0$),
and complementary slackness ($\lambda_i g_i = 0$). When you see
"closed-form solver" in the paper, it means KKT-stationarity has
been solved analytically for the inner WMMSE step at fixed dual
variables, and the outer dual ascent on $(\bm\lambda, \nu)$ is what
runs iteratively. It is exactly the structure of Lagrangian duality
you have probably seen in optimisation 101, just with quadratic forms
in $\mathbb{C}^M$.

\textbf{QCQP (quadratically constrained quadratic program).} A problem
with quadratic objective and quadratic constraints. The exposure-
constrained sum-rate problem is one of these (after the WMMSE
substitution). General QCQPs are NP-hard, but the structure you have
(PSD constraints, ridge regulariser, KKT closed form for the inner
step) reduces yours to a sequence of $M\times M$ Cholesky solves with
a Newton-like outer loop on the dual variables.

\textbf{Dual variables $\bm\lambda$, $\nu$.} One non-negative scalar
per constraint. $\lambda_u$ is the "shadow price" on body $u$'s
absorbed-power budget: how much sum-rate would gain if you were given
$1\,$\textmu W more headroom on body $u$. $\nu$ is the same thing for
the BS transmit-power budget. In the slack regime,
$\lambda_u = 0$ for all bodies, the exposure operators drop out of
\cref{eq:closedform-companion}, and the precoder reduces to plain WMMSE
zero-forcing. In the binding regime, the active set of $\lambda_u$ is
non-zero and the precoder steers \emph{against} the exposure modes of
constrained bodies. That is the entire control story.

\textbf{ECBF (exposure-constrained beamforming).} Generic name for any
precoder that solves a sum-rate / SINR problem under SAR or
absorbed-power caps. Hochwald 2014 introduced the SAR-matrix version
for a single body. Ying 2015 / 2017 generalised to multi-user. Your
paper generalises again to a heterogeneous body population with
soft per-body budgets.
\end{defs}

\begin{defs}[Architecture terms]
\textbf{ISAC (integrated sensing and communications).} Use the same
panel and waveform for radar-style sensing in addition to comms. In
this paper, ISAC means monostatic radar at the BS panel: send a
chirp, listen for the echo, detect bodies that have no phone uplink.
This is how the BS learns about tier-C bystanders.

\textbf{CPI (coherent processing interval).} The window over which
radar pulses are coherently integrated. $10\,$ms CPI gives a $10\,$dB
gain in SNR per decade if the target is stationary. The paper's link
budget shows that a $10\,$ms CPI is enough to detect a $0\,$dBsm body
at $50\,$m with $P_d = 0.99$ at $P_{fa} = 10^{-6}$.

\textbf{CSI (channel state information).} The BS's estimate of $\hh_k$
for each served user. Comes from uplink pilots (the UE transmits a known
sequence; the BS measures what it received). Used both to decide the
precoder for the next slot \emph{and}, in this paper, to calibrate the
ray-traced path amplitudes against reality.

\textbf{Path dictionary $\mathcal{D}$.} A precomputed library of ray
sets indexed by candidate body positions in the cell. Run the ray
tracer once over the static scene-and-array geometry on a grid of
$1\,$m spacing across the cell, store $(\khat_n, \bm\psi_n,
\alpha_n^{(d)}, j(n))$ per probe point. Then at runtime, looking up
the path set for a body at position $\rr_0$ is an $O(1)$ lookup, not
a fresh ray trace. Pose still varies (the body's mesh is different
when posed), but the path geometry on its $1\,$m grid cell does not.
This is what makes the loop run in real time.

\textbf{Translation phasor.} If the body moves by $\mathbf{t}$ inside
its grid cell (so the path \emph{geometry} doesn't change, only the
phase along each ray to the body's centre), the per-body Gram matrix
$\Mmat^{(u)}$ updates by a diagonal complex-exponential factor:
$\Mmat^{(u)}(\rr_0 + \mathbf{t}) = \Phimat(\mathbf{t})^H
\Mmat^{(u)}(\rr_0) \Phimat(\mathbf{t})$. This avoids re-running the
ray tracer for sub-grid translations. It is just the path-space
version of the obvious thing: shift the observation point, every
plane-wave path picks up an extra phase $\exp(-i k_0 \khat_n \cdot
\mathbf{t})$.

\textbf{Tiered telemetry.} Four classes of body, distinguished by what
the BS knows about them: A (served users, full GPS+IMU), B (cooperating
non-users, same telemetry but no served traffic), C (sensed bystanders,
position from ISAC, pose from worst-case envelope), D (undetected,
fall back to occupancy assumption). The reason this matters: you have
to constrain bystanders too, and bystanders by definition do not
report anything to the BS. The tier system makes that operational.
\end{defs}

\begin{defs}[The "regime" jargon]
\textbf{Slack regime.} The exposure budgets $L^{(u)}$ are larger than
the achievable absorbed power under any feasible precoder. Dual
variables $\lambda_u$ all sit at zero. Exposure constraints do not
bind. The precoder reduces to WMMSE / regularised ZF. This is what
the plaza scenario shows: $\Pabs/L_{\mathrm{RL}} \approx 10^{-3}$, far
below the cap.

\textbf{Binding regime.} At least one $\lambda_u > 0$. The exposure
constraints actively shape the precoder. The closed-form QCQP machinery
does real work. This regime activates at higher transmit power,
closer range, lower user counts, or under tighter caps (the post-2014
Brussels arrêté at $6\,$V/m, Italian outdoor attention values).

\textbf{Why the regime split matters.} Most of the paper's quantitative
results are in the slack regime, because that is what the canonical
plaza load actually is. The pose-information ablation needs the
binding regime to even produce a non-trivial result. The chronic-dose
argument and the spatial heatmap apply across both regimes.
\end{defs}

\section{The exposure operator and the precoder math, in your language}

You already know
\begin{equation}\label{eq:Q-recap}
  \Qmat^{(u)} = \int_{\Sigma^{(u)}} \tilde{\mathbf{G}}(\rr)^H
  \tilde{\mathbf{G}}(\rr) \, \mathrm{d} A
\end{equation}
from the summary paper. The new factorisation that the JSAC paper
uses is
\begin{equation}\label{eq:Q-factor-recap}
  \Qmat^{(u)} = \Jmat^T\,\Mmat^{(u)}\,\Jmat,
\end{equation}
which is just \eqref{eq:Q-recap} pushed through the path-space
identity $\tilde{\mathbf{G}} = \tilde{\bm\Psi} \Phimat \Jmat$. The
content of this is purely organisational, but it lets you separate
\textbf{what depends on the array} ($\Jmat$, fixed forever) from
\textbf{what depends on the body and pose} ($\Mmat^{(u)}$, refreshes at
$\sim 100\,$ms) from \textbf{what depends on body position}
($\Phimat$, refreshes at $\sim 1\,$s as the body walks). The
two-timescale separation in the paper rides this factorisation.

So far, this is just bookkeeping. The new content is what to do with
$\Qmat^{(u)}$ once you have one per body.

\subsection{From single-user SAR-aware to multi-user sum-rate}

Hochwald 2014 wrote the single-user version: maximise SNR
$|\hh^T \xx|^2 / \sigma_n^2$ subject to $\xx^H \mathbf{S} \xx \leq L$
and $\|\xx\|^2 \leq P$. It is a Rayleigh-quotient problem with one
extra ellipsoid constraint, and the closed form is
$\xx^\star \propto (\lambda \mathbf{S} + \nu \mathbf{I})^{-1} \hh^*$
(Ying 2015). One inverse, two scalar dual variables, done.

The JSAC paper goes multi-user. Now there are $K$ served users (each
wants their own data stream) and $K + B$ exposure constraints (one per
body, served users count, bystanders count too). The objective is
sum-rate $\sum_k \log(1 + \mathrm{SINR}_k)$, where SINR for user $k$
is $|\hh_k^T \ww_k|^2 / (\sigma_n^2 + \sum_{k' \neq k} |\hh_k^T
\ww_{k'}|^2)$. The interference-in-the-denominator is what makes
sum-rate non-convex.

\textbf{This is where WMMSE enters.} The trick is that
$\log(1 + \mathrm{SINR}_k)$ has the same global maxima (in $\ww_k$)
as a quadratic surrogate
$\log u_k - u_k e_k(\ww_k)$ for the right per-user weight $u_k$ and
MSE $e_k$. The proof is one paragraph of algebra (Christensen 2008).
What you get out of it is: at fixed $u_k, v_k$, the objective is
quadratic in $\ww_k$, and the exposure and power constraints are
quadratic too. You have a clean QCQP, and the inner KKT for $\ww_k$
gives
\begin{equation}\label{eq:closedform-companion}
  \ww_k^\star = \bigl(\Qmat_{\mathrm{tot}}(\bm\lambda) + \nu \mathbf{I}
  + \text{interference matrix}\bigr)^{-1}\,\gvec_k,
\end{equation}
where $\Qmat_{\mathrm{tot}}(\bm\lambda) = \sum_u \lambda_u \Qmat^{(u)}$
is the dual-weighted aggregate exposure operator and $\gvec_k =
u_k v_k \hh_k$ is the WMMSE drive.

\textbf{Why this is interesting.} The denominator
$(\Qmat_{\mathrm{tot}} + \nu \mathbf{I} + \text{interference})$ has
exactly three pieces:
\begin{itemize}[leftmargin=1.4em,itemsep=2pt]
  \item $\Qmat_{\mathrm{tot}}$: \emph{don't shine on bodies}
  \item $\sum_{k'} u_{k'} |v_{k'}|^2 \hh_{k'} \hh_{k'}^H$:
  \emph{don't interfere with other users}
  \item $\nu \mathbf{I}$: \emph{don't blow up the transmit power}
\end{itemize}
That is regularised zero-forcing with one extra term: the body-
aggregated exposure operator. You are inverting the same matrix shape
as ZF, but with an exposure-aware regulariser inside.

\subsection{The rank-one limit and why ZF is "the rank-one case"}

Set $\Qmat^{(u)} = \sigma_u \mathbf{a}_u \mathbf{a}_u^H$ for every body
(every body collapses to a scalar steering vector). Stack
$\mathbf{a}_u$ into a matrix $\mathbf{A}$. Then
$\Qmat_{\mathrm{tot}} = \mathbf{A}^H \mathbf{D} \mathbf{A}$
with $\mathbf{D} = \mathrm{diag}(\lambda_u \sigma_u)$. The closed form
\eqref{eq:closedform-companion} reduces to
$\ww_k^\star \propto (\mathbf{A}^H \mathbf{D} \mathbf{A} + \nu
\mathbf{I})^{-1} \hh_k$, which is the regularised inverse with
\emph{user channels and bystander steering vectors stacked together}.
That is exactly Peel-Hochwald-Swindlehurst regularised ZF, applied to
the union of served-user channels and bystander steering directions.

So the message of the rank-one limit is that the paper's proposal does
not invent a new precoder. It shows that the multi-body exposure-
constrained problem and ZF are the \emph{same} machine, with full
$\Qmat^{(u)}$ replacing rank-one steering vectors. Soft budgets cost
less than hard nulls because they share DoF instead of consuming one
per body.

\begin{nb}[The "soft budget" part that confused you]
With hard nulls (ZF), each body or user costs exactly one precoder
DoF, and you saturate at $K + B \approx M$. With soft budgets, the
$\lambda_u$ trade off how aggressively each body's exposure is
suppressed against the total transmit power. Two nearby bodies that
share an angular subspace can be jointly soft-budgeted using less than
two DoF. With $K + B = 50$ and $M = 64$, hard nulling exhausts the
array, but soft budgeting (with $\Qmat^{(u)}$ rank 3 to 4 per body)
is well within the closed-form regime.
\end{nb}

\subsection{What the dual ascent actually does}

The closed form \eqref{eq:closedform-companion} solves the inner
WMMSE problem at fixed $\bm\lambda, \nu$. The outer loop:
\begin{enumerate}[leftmargin=1.6em,itemsep=2pt]
  \item Initialise $\bm\lambda = 0$, $\nu$ from the power budget.
  \item Solve \eqref{eq:closedform-companion} for $\{\ww_k^\star\}$.
  \item Evaluate constraint residuals $F_u = \sum_k \ww_k^{\star H}
  \Qmat^{(u)} \ww_k^\star - L^{(u)}$.
  \item Newton step on $\bm\lambda$: bodies with $F_u > 0$ get
  $\lambda_u$ increased, bodies with $F_u < 0$ stay at zero.
  \item Repeat until residuals are below tolerance.
\end{enumerate}
This converges in 10-30 iterations cold-started, $\sim 5$ warm-started
across slots. Each iteration is one $M \times M$ Cholesky solve at
$O(M^3) = O(64^3)$. The published median is $31\,$ms on CPU for
$K=B=25$, $M=64$. GPU vmap brings that under $10\,$ms.

That is the whole solver. It is dual-decomposition for QCQP, nothing
exotic. The exotic part is that the constraints are exposure operators
on bodies, computed from your TAP-paper machinery.

\section{The closed-loop architecture, ingredient by ingredient}

The paper's §IV-V is about \emph{how} the BS actually maintains
$\Qmat^{(u)}$ for every body in real time. This is the part you wrote
the spine doc to clarify. Here is the cleaner version.

\subsection{Why a path dictionary}

You cannot afford to ray-trace per body per slot. A naive cost is
$50\,\text{bodies} \times 1\,\text{ms slot} = 50\,000\,$ray traces per
second. Even on GPU that is not happening.

The dictionary $\mathcal{D}$ is precomputed offline, once per scene-and-
panel installation. Run the ray tracer over the static geometry
(BS panel + buildings + ground, no bodies) on a $1\,$m grid of
candidate body-centre positions. Store the resulting path set
$\{\khat_n, \bm\psi_n, \alpha_n^{(d)}, j(n)\}$ per grid cell. For a
$100 \times 100\,$m cell at $1\,$m spacing with $\leq 20$ paths per
cell, this fits in tens of megabytes.

At runtime, when body $u$ is at position $\rr_0$, its path set is
$\mathcal{D}[\rr_0]$, an $O(1)$ lookup. Pose changes only the per-body
Gram $\Mmat^{(u)}(\rr_0, \bm\theta_u)$, which you recompute from the
SMPL-X mesh. Position changes inside the grid cell are handled by the
translation phasor.

\textbf{Why this works.} The ray geometry between BS and a $1\,$m
grid cell does not change much when you nudge the body $0.3\,$m
inside the cell. The phase along each ray does change, hence the
phasor. Pose changes the body mesh, so the integrand
$\tilde{\bm\psi}_n^H \tilde{\bm\psi}_{n'}$ over the body surface
changes, hence the per-pose $\Mmat^{(u)}$ recomputation. But the
ray tracer is amortised over the offline dictionary build.

\subsection{Why CSI calibration}

The path \emph{directions} from the dictionary are accurate (ray
tracing on real building geometry is fine). The path \emph{amplitudes}
are not. Real deployments suffer from antenna calibration error,
material mismatches against the OSM database, missing facade detail,
diffuse scattering from window glass that the ray tracer
discretised wrong. So you fit a complex per-path scale $\beta_n$
against the served users' uplink CSI:
\[
  \hat{\bm\beta}_k = \arg\min_{\bm\beta} \|\hh_k^{\mathrm{meas}} -
  \mathrm{diag}(\bm\beta) \cdot \hat{\hh}_k\|^2 + \rho \|\bm\beta -
  \mathbf{1}\|^2.
\]
The ridge term biases toward the in-silico amplitudes when the LS
problem is underdetermined or noisy. The fitted $\bm\beta$ propagates
to every body that shares dictionary entries with user $k$, since the
dictionary entries are shared across bodies in nearby grid cells.
This is what closes the in-silico to deployed gap. The residual is
characterised in \cref{fig:csi-companion}.

\begin{figure}[!h]
  \centering
  \includegraphics[width=0.9\linewidth]{residual_vs_snr.png}
  \caption{CSI calibration residual vs uplink SNR for $M=64$, three
  path counts, plain LS vs ridge LS. The dashed black baseline at
  $-4.7\,$dB is the no-calibration error (set by the perturbation
  magnitude). \textbf{Plain LS (solid)} only beats no-calibration when
  $N \leq M/2$ \emph{and} SNR is high; at $N=60 \approx M$ it amplifies
  noise badly. \textbf{Ridge LS (dashed)} converges at all path counts
  and gives a clean trend with SNR. Operational rule: keep $N \leq
  M$, cluster correlated paths if needed, and always use the ridge.}
  \label{fig:csi-companion}
\end{figure}

\subsection{The tier system}

Bodies in a cell come in four flavours, and the paper has to handle
all four because the budget constraint applies to all of them:

\begin{table}[!h]
\centering
\small
\begin{tabular}{@{}p{0.05\linewidth}p{0.22\linewidth}p{0.32\linewidth}p{0.28\linewidth}@{}}
\toprule
Tier & What they are & Position from & Pose from \\
\midrule
A & Served users & 3GPP UL pilots + GPS via MNO app & Phone IMU $\to$ SMPL-X via Madgwick or VQF (existing literature) \\
B & Cooperating non-users (registered phone, no traffic) & Same as A & Same as A \\
C & Sensed bystanders (no phone, no MNO context) & ISAC monostatic radar at the panel & Worst-case Cauchy envelope (no pose telemetry) \\
D & Undetected & Occupancy grid (regulator-defined) & Worst-case envelope \\
\bottomrule
\end{tabular}
\caption{The tier table from the paper, annotated. The reason this
exists is that bystanders are constraints too, and bystanders do not
report to the BS by definition.}
\end{table}

The thing that took you a while to internalise: \textbf{the BS does
not need the body to tell it the angle of attack.} The BS already
runs its own ray tracer. What it needs from the body is
\emph{position}, full stop. Pose is a second-order refinement that
shapes $\Qmat^{(u)}$'s trace and leading eigenvector. For tier A and
B you get pose from the phone IMU. For tier C, you fall back to the
Cauchy worst case ($D \leq 2$ on convex bodies, projected area at
its peak orientation). For tier D, you assume an occupancy envelope.

\subsection{The cadence table, decoded}

\begin{table}[!h]
\centering
\small
\begin{tabular}{@{}lll@{}}
\toprule
Quantity & Cadence & Per-update cost \\
\midrule
UE CSI $\hh_k$ & $1\,$ms & $O(M)$ from PHY pilots \\
Precoder $\ww_k^\star$ & $1\,$ms & $31\,$ms median solve, $K=B=25$, $M=64$ \\
Translation phasor & $\sim 1\,$s & $O(N)$ per body \\
Pose-driven $\Mmat^{(u)}$ & $\sim 100\,$ms & $7.5\,$ms per body (SMPL-X) \\
Path dictionary $\mathcal{D}$ & static & offline, minutes \\
CSI calibration $\bm\beta_k$ & $\sim 1\,$s & LS fit, $O(NM)$ \\
\bottomrule
\end{tabular}
\caption{The cadence ladder. Slow things at the top, fast things at
the bottom. Each row's cost is amortised over its cadence: the
precoder solve is $31\,$ms per slot, but warm-starting across slots
brings it down. The pose refresh dominates per-body GPU cost.}
\end{table}

\textbf{Why "$1\,$ms slot but $31\,$ms solve" is fine.} The slot
cadence is what the PHY transmits at; the solve runs in pipelined
fashion across slots and is amortised. In production, you warm-start
$\bm\lambda^{(t)}$ from $\bm\lambda^{(t-1)}$ and the dual converges
in five iterations. Real BS implementations do this kind of pipelined
solve constantly, this is not novel.

\section{What the experiments are actually showing}

The plaza experiment is the heart of the paper. There are seven figures
and they each tell one story. Let me walk you through them with the
images inline.

\subsection{The plaza geometry}

\begin{figure}[!h]
  \centering
  \includegraphics[width=0.65\linewidth]{plaza_top.png}
  \caption{Brussels Grand Place geometry as the experiment sees it.
  The orange square at the bottom is the BS panel on the south facade
  at $8\,$m height with $10^\circ$ downtilt. The labelled red dots are
  OSM-verified entry points where bodies enter and exit the plaza on
  origin-destination flux traces. Cell extent is roughly $80 \times
  80\,$m. Bodies are $5$ to $55\,$m from the panel.}
\end{figure}

You probably already had a mental picture, but it helps to see it
explicitly. The bodies walk between these seven entry nodes on AMASS
walk-cycle traces. The BS sees them as a cloud of moving SMPL-X
phantoms, ray-traces to each at $\sim 100\,$ms cadence, and updates
$\Qmat^{(u)}$ accordingly.

\subsection{Operator rank: 3 to 4, basically always}

\begin{figure}[!h]
  \centering
  \includegraphics[width=0.95\linewidth]{rank_cdf.png}
  \caption{Empirical eigenvalue spectrum of $\Qmat^{(u)}$ over 20
  Thelonious bodies in a $20$-$80\,$m ring around the panel. The
  $\lambda_k / \lambda_1$ y-axis is logarithmic and falls off a
  cliff at $k = 3$ (specular) or $k \approx 5$ (UMa-LOS). The
  effective rank is the index where the trace fraction drops below
  some threshold $\epsilon$.}
\end{figure}

This is the easiest figure in the paper. The story:

\begin{itemize}[leftmargin=1.4em,itemsep=2pt]
  \item In the plaza-specular path model (LOS + ground + 2 facades),
  $\Qmat^{(u)}$ has effective rank 3 at $\epsilon = 10^{-2}$, max 4 at
  $\epsilon = 10^{-3}$.
  \item In the stochastic 3GPP UMa-LOS model (60 subpaths per cluster),
  the rank is even smaller at the leading edge: median 1 at
  $\epsilon = 10^{-1}$, climbing to 3 at $10^{-3}$.
\end{itemize}

\textbf{Why so low?} Because the $8 \times 8$ panel resolves $\sim
14^\circ$ in azimuth at $26\,$GHz. A body subtends maybe two beams
across the torso. The dominant body-coupling subspace collapses into
3 to 4 distinguishable BS-side beam directions, regardless of how
much multipath the path model gives you. \textbf{Adding more
multipath raises the trace, not the rank.} This is a panel-side
affordance, not a model simplification.

\textbf{Why this matters for the algorithm.} The closed-form solver
inverts $\Qmat_{\mathrm{tot}}(\bm\lambda) + \nu \mathbf{I}$, which is
a $64 \times 64$ matrix. When each $\Qmat^{(u)}$ has rank 3 to 4 and
you have $K+B = 50$ bodies, $\Qmat_{\mathrm{tot}}$ has at most
$\sim 4 \times 50 = 200$ effective directions, well above $M = 64$, so
the soft-budget formulation is well-posed. Hard nulling 50 bodies
into a 64-dim array would be hopeless; the rank-few structure is what
saves it.

\subsection{The hero figure (which is a null result, and why that's OK)}

\begin{figure}[!h]
  \centering
  \includegraphics[width=\linewidth]{hero_pareto.png}
  \caption{(a) ECDF of $\Pabs / L_{\mathrm{RL}}$ over 450{,}000
  body-slots in the canonical plaza load. The violation zone is at
  $x = 1$ (right margin). The data clusters at $x \approx 10^{-3}$.
  All five precoders satisfy the cap by $\sim 1000\times$. (b) Mean
  sum-rate per precoder. ZF wins on rate in plaza-specular paths
  (44\,Gbps because the narrow-angle path set keeps $\hh \hh^H$
  well-conditioned for $K=25$ on $M=64$), but degenerates in UMa-LOS
  where the effective channel rank falls below $K$.}
\end{figure}

\textbf{What you wanted this to show.} Some baseline (MRT or ZF)
violates the cap and your ECBF does not. A clean Pareto with your
precoder dominating.

\textbf{What it actually shows.} The cap is slack by $1000\times$. All
five precoders sit on the same CDF on the budget side. ZF separates
on the rate side (because it nulls inter-user interference), but on
the compliance metric, no one violates and no one is interesting.

\textbf{Why the paper survives this null.} You pivoted (correctly) to
a regime-characterisation paper. The hero figure now does three jobs:
(i) it shows the algorithm is correct (it satisfies the cap), (ii)
it locates the binding regime precisely (this scenario is not it),
(iii) it motivates the chronic-dose argument as the actual
deployment value-prop.

The handoff doc says "the plaza data is a null on the instantaneous
metric." Yes. The pivot is to chronic-dose ECDFs and the spatial
heatmap, both of which separate even when the instantaneous compliance
is easy.

\subsection{Chronic dose ECDFs by tier}

\begin{figure}[!h]
  \centering
  \includegraphics[width=0.85\linewidth]{chronic_dose.png}
  \caption{ECDF of body-integrated absorbed energy over a $300\,$s
  origin-destination flux trace, stratified by tier.
  Tier A served users $n=25$, B cooperating $n=15$, C bystander
  $n=10$. Median: A $\sim 12.6\,$mJ, B $\sim 8.4\,$mJ, C $\sim 9.1\,$mJ.}
\end{figure}

\textbf{The story.} Even when instantaneous compliance is easy,
\emph{chronic dose} (integrated $\Pabs$ over a 5-min walk) separates by
tier. Served users (red) carry the most because their precoders steer
toward them. Cooperating non-users (blue) carry the least because they
get only side-lobe spillover and the precoder is not aimed at them.
Bystanders (black) sit in between because their dose is also side-lobe
dominated, but they are not steered \emph{away from} either.

\textbf{The 1.5$\times$ ratio between A and B is the practical value
proposition.} An operator that wants to argue chronic-dose reduction
for the bystander population can point at this and say "even at slack
instantaneous compliance, our DL serves users at 1.5$\times$ the dose
we deposit on bystanders." This is the deployable angle when the
instantaneous cap is not binding.

\textbf{Caveat the paper flags.} Earlier shorter (20\,s) traces
over-stated the gap at $3\times$. Longer traces converge to $1.5\times$
because side lobes dominate cumulative dose, not the steered beams.
You should be careful when reading the paper not to claim "$3\times$"
in conversation: that was a stale number from a static-position trace.

\subsection{Spatial heatmaps in slack vs binding regimes}

\begin{figure}[!h]
  \centering
  \includegraphics[width=\linewidth]{spatial_heatmap.png}
  \caption{(a) Slack regime, multi-body ECBF: forward-sector hotspot
  $\sim 25\,$m north of the panel, peaks at $\sim 4.8 \times 10^{-4}
  L_{\mathrm{RL}}$. (b) Binding regime ($43\,$dBm Tx, $E_{\mathrm{RL}}
  = 6\,$V/m post-2014 Brussels / Italy), ZF: cell-mean median
  $0.06\,L_{\mathrm{RL}}$, hot cells $0.4\,L_{\mathrm{RL}}$, but per-
  body-slot tail to $3.4\,L_{\mathrm{RL}}$ (i.e. ZF violates).}
\end{figure}

\textbf{Two stories in one figure:}

\textit{Slack regime (left).} The deposition map looks like the array's
projected gain pattern. Hot zones are where the broadside spills the
most energy. Bodies that walk through these zones absorb more. Nothing
is close to the cap.

\textit{Binding regime (right).} Higher Tx power, tighter RL ($6\,$V/m
not $14.57$). Cell-averaged compliance looks acceptable
($0.06\,L_{\mathrm{RL}}$ median), but the per-body tail violates: rare
body-slots reach $3.4\,L_{\mathrm{RL}}$ where a body walks through a
high-gain ridge between two served-user beams. \textbf{ZF nulling is
geometric, not exposure-aware.} Your ECBF holds the per-body
$99^{\mathrm{th}}$ percentile at $1.0\,L_{\mathrm{RL}}$ by
construction.

\textbf{This is the key message.} The case for budget-aware precoding
strengthens with array size. At $M = 256$ the ZF violation rate climbs
from $0.10\%$ (M=64) to $18.7\%$, even with 231 nulling DoFs over 25
served users. A larger array focuses energy more tightly, and
geometric ZF nulls do not see exposure tails. The operational argument
holds even when the cell-averaged metric looks fine.

\subsection{Pose-information ablation}

\begin{figure}[!h]
  \centering
  \includegraphics[width=\linewidth]{pose_info_gain.png}
  \caption{Pose ablation in the binding regime. (a) Slot-level $\Delta$
  sum-rate (aware $-$ ablate) ECDF: bimodal, with a large left tail
  (-10\,Gbps) and a smaller right tail (+20\,Gbps). MRT and WC back-off
  collapse to $\Delta = 0$ because they don't consume $\Mmat^{(u)}$.
  Oracle is also $\Delta = 0$ because it always uses posed-Q. (b) Per-
  body $P_{\mathrm{abs}}$ ECDF: pose-aware (blue) dominates pose-
  ablated (red dashed) at every quantile, with the red curve hitting
  $L_{\mathrm{RL}}$ much sooner.}
\end{figure}

This figure is more nuanced than the others, and you should re-read the
paper's §V.E with the following picture in mind:

\textbf{The compliance side (right panel).} Pose-aware solver has
$0.96\%$ violation rate vs $4.58\%$ without pose. \textbf{$4.7\times$
compliance gain.} This is the win.

\textbf{The sum-rate side (left panel).} Bimodal. Sometimes pose-aware
is much \emph{worse} on sum-rate ($35.7\%$ of slots have $\Delta < -1\,$Gbps,
left tail), sometimes much \emph{better} ($12.6\%$ of slots have
$\Delta > +1\,$Gbps, right tail). The negative tail is solver-side: the
posed operator is \emph{tighter} than the Cauchy envelope, so the
feasible cone is smaller, the dual-Newton sweep is harder to converge
in 8 iterations, and the solver falls into min-absorption fallback
more often ($83.6\%$ pose-aware vs $55.1\%$ Cauchy).

\textbf{What this means.} Pose information buys a \emph{tighter
exposure profile} but costs solver iterations. The trade-off vanishes
once you uncap the Newton sweep and replace the Cauchy envelope with
a tighter Lanczos-projected operator (the codebase already has the
ingredients). The oracle (which uses posed-Q always with no fallback)
achieves $11.79\,$Gbps with full compliance, showing the upper bound
of what the solver could reach.

\textbf{Don't read this figure as "pose information is a wash."} It is
not. The compliance gain is real and large. The sum-rate cost is a
solver-engineering issue, not a fundamental limit.

\subsection{Tier-C ISAC link budget}

\begin{figure}[!h]
  \centering
  \includegraphics[width=\linewidth]{detection_vs_range.png}
  \caption{Monostatic ISAC link budget at the panel. (a) Integrated
  SNR vs range for three CPI durations. Plaza range ($\sim 50\,$m)
  hits $42\,$dB at $10\,$ms CPI. (b) Detection probability (Albersheim,
  Swerling-0) saturates by $50\,$m and reaches $P_d = 0.99$ at
  $137\,$m for $10\,$ms CPI.}
\end{figure}

\textbf{What this is showing.} Monostatic radar detection of bodies
(treated as $0\,$dBsm RCS targets) at the BS panel, using the same
$8\!\times\!8$ array at $26\,$GHz. The point is that the link budget
\emph{closes}: bystander detection is feasible in the plaza geometry
under standard radar assumptions. $137\,$m at $P_d = 0.99$ with
$P_{fa} = 10^{-6}$ is a comfortable margin over plaza range.

\textbf{What this is not showing.} Tracking, clutter rejection,
spatial discrimination of multiple bodies, vital-signs at plaza range.
Those are out of scope (and the brainstorm rounds explicitly walked
back the vital-signs claim). This figure says "yes, you can detect
that something body-shaped is at this range and angle." The Cauchy
envelope handles everything else.

\textbf{Why this is one figure not seven.} Round 3 of the brainstorm
killed the ISAC heartbeat / respiration angle (no plaza-range
literature exists for that). What's left is presence + angular
localisation, which closes by a wide margin and lives in one
sub-figure of one section.

\subsection{Cauchy bound failure}

\begin{figure}[!h]
  \centering
  \includegraphics[width=\linewidth]{tightness.png}
  \caption{Tightness of the rank-1 Cauchy operator bound (Theorem 2 of
  paper\_v2). The $x$-axis is $10 \log_{10}(\xx^H \Qmat \xx /
  \text{RHS})$. Negative means the bound holds, positive means it is
  violated. Different curves are different precoder ensembles; the
  bound applies to a precoder direction. \textbf{MRT to LOS direction
  (red, leftmost):} bound holds by 21 to $27\,$dB. \textbf{Random,
  ECBF, ZF (blue/orange/green/purple):} bound \emph{breached} at
  the median of the distribution.}
\end{figure}

This is the figure that forced the biggest narrative correction in the
JSAC paper. Read this carefully because it changed the paper's claims.

\textbf{What you wanted Cauchy to do.} Provide a universal upper bound
on $\xx^H \Qmat^{(u)} \xx$ in terms of the body's projected area and
the BS-side directivity, so the twin could satisfy the budget without
pose telemetry by applying the Cauchy worst case.

\textbf{What Cauchy actually does.} The bound holds for $\xx$ that
points in the BS-side steering direction $\mathbf{a}_{\mathrm{BS}}^{(u)}$.
For an MRT-LOS precoder it is loose by 21 to $27\,$dB. For random and
ECBF precoders it \emph{breaches} 49 to $70\%$ of the time, with
median ratios up to $+6.34\,$dB above the bound.

\textbf{Why it breaks.} The leading exposure mode $\mathbf{v}_1$ of
$\Qmat^{(u)}$ does not align with the BS-side steering vector
$\mathbf{a}_{\mathrm{BS}}$. The inner product
$|\langle \mathbf{v}_1, \mathbf{a}_{\mathrm{BS}}\rangle|^2$ sits in
$[0, 0.19]$ across all evaluated jobs. So the projected-area bound,
derived under the assumption that the absorbed power is concentrated
along the BS-pointed direction, simply misses where the absorption
actually peaks. The projected-area integral averages $D \leq 2$ over
the sphere, but the body's leading exposure mode picks out a specific
direction not on the BS axis.

\textbf{The fix.} Theorem~2 of paper\_v2 was demoted from a universal
upper bound to a \emph{conditional} bound (Appendix C of the JSAC
paper). The conditional form holds when the precoder is restricted to
the conic hull of a chosen direction set $\mathcal{A}$ that includes
$\mathbf{a}_{\mathrm{BS}}$. The price is that you have to commit to
the direction set up front, and any precoder outside that set falls
back to direct $\xx^H \Qmat^{(u)} \xx$ evaluation (which is fine,
just expensive in the binding regime).

\textbf{Why this is interesting and not just an embarrassment.} It is
a clean geometric fact about how 64-element panels and human bodies
couple. The body's dominant absorption mode is shaped by the body's
geometry under whatever angular spectrum hits it; the BS steering
vector points at the user's location. These are not the same direction
unless the user happens to be the body and the body happens to face
the BS at exactly the right orientation. So the disconnect is
generic, and it is now a paragraph in the discussion section that
flags this for any future compliance-bound construction. The honest
version is the right version.

\section{What's actually new vs what you already had}

The paper has six contributions, and four of them existed in pieces
before this paper. Here is the honest accounting.

\begin{table}[!h]
\centering
\small
\begin{tabular}{@{}p{0.30\linewidth}p{0.62\linewidth}@{}}
\toprule
Contribution & State of the art before this paper \\
\midrule
Coherent $\Qmat^{(u)}$ from path set & In your monograph and the JSAC paper's §II. The path-space factorisation $\Jmat^T \Mmat \Jmat$ is from your summary paper. Reused, not new. \\
Single-body SAR-aware precoding & Hochwald 2014, Ying 2015/17. Existed. Single-body, single-pose, FDTD-calibrated $\mathbf{S}$. \\
Multi-user, multi-body QCQP closed form & \textbf{This is new.} Generalisation of Ying-2015 to a heterogeneous population, with bystanders as constraints. The KKT structure and the rank-one ZF limit are new statements. \\
Scene path dictionary $+$ translation phasor $+$ CSI calibration & \textbf{This is new as an architecture.} Each piece individually is standard (ray tracing for scene; channel charting for calibration; path-space phase shift is trivial), but the composition that lets the loop run on real BS cadences is the actual operational contribution. \\
Tiered telemetry (A/B/C/D) & \textbf{New as a framing.} The pieces (3GPP location services, IMU mocap, ISAC at panel, occupancy envelopes) are individually deployed elsewhere. The categorisation that makes bystander compliance operational is yours. \\
Plaza-scale empirical regime characterisation & \textbf{This is new and is the paper's empirical anchor.} No prior work has run a 50-body multi-user ECBF at plaza scale on real OSM geometry under a real RL. The rank-3-to-4 result, the $7.5\,$ms/body GPU number, the $4\times 10^4$ slack margin, the chronic-dose tier ECDFs, and the Cauchy bound failure are all empirical contributions specific to this paper. \\
\bottomrule
\end{tabular}
\caption{What the paper adds to the literature, and where it is
standing on prior work.}
\end{table}

So when you re-read the paper and feel "is this novel enough?", the
honest answer is: it has one new theoretical contribution (multi-body
QCQP closed form including bystanders), one new architectural
contribution (the dictionary + tiers + cadence ladder), and a real
empirical anchor at plaza scale that no one has done before. The
single-body machinery and the body-side dosimetry are reused from
your earlier work.

\section{Where the paper is rough on purpose}

The handoff doc lists three places. Read them with the following
mental model: each is a known limitation that does not affect the
qualitative regime story.

\begin{itemize}[leftmargin=1.4em,itemsep=2pt]
  \item \textbf{Single seed, $20\,$s window, Shannon proxy.} A multi-
  seed average would smooth the bystander tail (currently $n=10$
  bodies, visibly coarse on the chronic-dose ECDF). A real Sionna NR
  PHY pass would shift absolute sum-rate numbers but should not affect
  the precoder collapse. Both are noted in the discussion.
  \item \textbf{Cauchy proof in Appendix~C}, appeals to the angular-
  spectrum decomposition without re-deriving the projected-area
  identity. The identity is classical and in your TAP paper.
  \item \textbf{Brussels averaging window} cited as $6\,$min per
  current practice; the 2024 ordinance number is a placeholder
  citation. Italian $15\,$V/m attention value is flagged TODO.
\end{itemize}

\section{What you can hand to the next paper}

The brainstorm rounds named several follow-ups that got sidelined.
Three of them are real next-paper material:

\begin{enumerate}[leftmargin=1.6em,itemsep=2pt]
  \item \textbf{Pose-differentiable $\Qmat^{(u)}(\bm\theta)$.} Round-1
  hero, deferred. SMPL-X is differentiable, your TAP machinery is
  closed form, so $\Qmat^{(u)}$ is differentiable in $\bm\theta$ in
  principle. A pose-loop SDP or proximal solver that directly
  optimises pose for compliance is genuinely new. Killed in this paper
  because the slack regime made it moot.

  \item \textbf{Multi-operator cooperative cap allocation.} Brussels
  cumulative cap is shared across operators in the same band. Each
  operator sees a slot of the cumulative budget. Mechanism design
  (Shapley, VCG) and the cooperation gap are open. You sidelined this
  in round 2.

  \item \textbf{Reflective intelligent surface in the binding-regime
  corner.} A one-line modification to $\Mmat^{(u)}$ in
  $\Qmat = \Jmat^T \Mmat \Jmat$. The RIS adds a controllable subspace
  to $\Mmat$. The open question is what RIS configuration buys back
  the most cap headroom for a given body distribution. Not in this
  paper, listed as future work only.
\end{enumerate}

\section{The thing to remember}

\textbf{The paper is regime characterisation.} It is not "we built a
new precoder that beats ZF". The new precoder \emph{contains} ZF as a
limit. The paper's claim is:

\begin{quote}
Under canonical mmWave plaza conditions on the Brussels Grand Place,
the Brussels reference level is slack by four orders of magnitude;
the binding regime is exclusion-zone geometry rather than coverage
geometry; the chronic-dose argument carries the deployable value at
plaza scale; pose information shapes the operator and gives a real
compliance lever in the binding regime; the projected-area Cauchy
bound is conditional, not universal, with the geometric mismatch
between the leading exposure mode and the BS-side steering vector
explaining the breach. The body-twin architecture is closed-form,
real-time-feasible at $7.5\,$ms/body GPU pose-refresh and $31\,$ms
median multi-body QCQP solve, and is composable from quantities the
network already maintains.
\end{quote}

That paragraph is the thing the JSAC reviewer needs to come away with.
Everything else is justification. Re-read the paper through this lens
and the structure makes sense: §II builds the operator (your existing
work, slimmed), §III states the multi-body QCQP and its closed form
(the one new theoretical thing), §IV-V build the architecture (the
new operational thing), §VI-VII run the plaza experiment and tell the
regime story, §VIII discusses, §IX concludes. The figures map onto
this structure one to one.

\bigskip\hrule\bigskip

\section*{If you forget everything else, remember these five things}

\begin{enumerate}[leftmargin=1.6em,itemsep=4pt]
  \item \textbf{The closed form $\ww_k^\star = (\Qmat_{\mathrm{tot}} +
  \nu \mathbf{I})^{-1} \gvec_k$ is regularised ZF with an exposure-
  aware regulariser.} It contains MRT, ZF, regularised ZF, and worst-
  case back-off as limits. Reduces to ZF when every $\Qmat^{(u)}$ is
  rank one.
  \item \textbf{The empirical $\Qmat^{(u)}$ rank is 3 to 4.} Set by
  panel angular resolution, not by multipath richness. This is what
  makes soft budgets cheap and the closed-form solver well-behaved at
  $K + B = 50$ on $M = 64$.
  \item \textbf{The plaza scenario is in the slack regime.} The
  Brussels RL is slack by $\sim 10^4$ at the canonical load. The
  binding regime activates inside $\sim 3\,$m of the panel at $K = 25$,
  or under tighter RL ($6\,$V/m post-2014 Brussels / Italy), or at
  larger arrays ($M = 256$ on the same scenario gives ZF violation
  rates of $18.7\%$).
  \item \textbf{Chronic-dose ECDFs separate by tier even in the slack
  regime.} Cooperating non-users absorb $\sim 0.4\times$ the served-
  user median over a 5-min walk. This is the deployable value-prop
  outside the binding ring.
  \item \textbf{The projected-area Cauchy bound is conditional, not
  universal.} It is loose by 21 to $27\,$dB on the LOS steering
  direction and breached by 49 to $70\%$ of unconstrained precoders.
  The conditional form (over a chosen direction set) is the right
  object to use, and the mechanism (leading exposure mode does not
  align with BS-side steering) is a clean geometric fact about how
  bodies couple to 64-element panels.
\end{enumerate}

\bigskip
\textit{End of companion. This document is for your own reading. Do
not cite it anywhere. The actual paper is paper\_jsac.tex; this is
just the friend who reads with you.}

\end{document}
